
\begin{frame}[t]{The Signal a Filing Carries}
\hypertarget{src:signal}{}
\small
A filing tells voters two different things, with very different reliability. As
an allegation about a candidate it is weak evidence, because a genuine charge
and a strategic one look alike from outside. As a fact about the race it is
strong evidence, because the filer is optimizing and sues the leading opponent
\ns{\citep{nakaguma2025}}. Litigation is precise about position and imprecise
about quality.
\vspace{4pt}

The precise half is what moves votes. A sharper read on who is ahead makes the
posterior gap between the top two less dispersed, so a reversal gets less
likely and the expected margin widens. That is the channel the estimates pick
up: the winner's share of the valid vote moves \WinShareCoef\ and the
runner-up's \RunnerUpCoef, widening the top-two margin by \MarginCoef\
($p=\MarginP$). A second leg is imported rather than estimated and points the
same way --- with the rank order common knowledge, supporters of candidates
ranked third and below desert toward the leader
\ns{\citep{myerson1993,anagol2016}} --- but these estimates do not separate it,
because the movement they show is between the top two. Both faces of the
previous frame produce a narrowing of this kind.
\vspace{4pt}

Two channels are closed by institutional design rather than by the estimates
happening to be small. Candidate registration closes in August and adversarial
filings run through October, so the field cannot respond ($p=\CandExecP$). And
voting is compulsory, so a voter leaving the contest adjusts the mark rather
than the attendance ($p=\TurnoutP$) --- which is what puts the response on the
blank and null ballot.
\vfill
\centerline{\hyperlink{app:theory}{\beamergotobutton{The model}}\quad\hyperlink{app:transferfoil}{\beamergotobutton{Why not a transfer model}}}
\end{frame}
